\section{Complete trace-case table}\label{app:cases}

Let the bounded source section on a prescribed line be either empty, a
point, or a compact segment, and let the ambient section be $[c,d]$.  The
construction retains the source section before clipping.  The possibilities
are summarized below; endpoint coincidences for several neurons are literal
geometric coincidences and therefore remain simultaneous.

\begin{center}
\begin{tabular}{@{}p{.21\textwidth}p{.20\textwidth}p{.22\textwidth}p{.25\textwidth}@{}}
\toprule
Bounded source data & Uncut strict trace & Relative strict trace & Retained finite datum\\
\midrule
$L\cap\overline U=\varnothing$ & $\varnothing$ & $\varnothing$ & none\\[.25em]
$L\cap\overline U=\{a\}$ & $\varnothing$ & $\varnothing$ & the closed point $a$\\[.25em]
$L\cap\overline U=[a,b]$, $L\cap U=\varnothing$ & $\varnothing$ & $\varnothing$ & the extreme points $a,b$\\[.25em]
$L\cap U=(a,b)$ & $(a,b)$ & $(a,b)\cap[c,d]$ & $a,b$ and two off-line frame points\\
\bottomrule
\end{tabular}
\end{center}

The set $(a,b)\cap[c,d]$ may be open, half-open, closed, a point, or empty.
The closed relative trace is always
\[
 [a,b]\cap[c,d].
\]
For the source-open-empty segment case, the entire closed segment is a
supporting face and hence contains no ambient-interior point of the output
polygon.  This is distinct from a lower-dimensional neural-code atom, which
may arise as the set-theoretic difference of several open fields even though
each individual field is full dimensional.

\paragraph{Whole-line source traces.}
The auxiliary bounded disk is introduced before the trace package is formed.
Thus an originally unbounded intersection with $L$ becomes a compact section
whose endpoints lie outside $F$.  Freezing that bounded section preserves the
entire clipped trace in $F$, without assigning finite endpoints to the
original unbounded line section.
